% !TEX root = ../manuscript.tex

\begin{description}
  \item[\textbf{Clustering}] Méthode d'apprentissage non supervisé visant à regrouper des observations similaires sans variable cible prédéfinie.

  \item[\textbf{Cross-validation (validation croisée)}] Technique d'évaluation qui consiste à diviser les données en $k$ sous-ensembles et à entraîner/tester le modèle $k$ fois.

  \item[\textbf{Data leakage}] Fuite d'information du jeu de test vers le jeu d'entraînement, entraînant une surestimation artificielle des performances.

  \item[\textbf{Feature engineering}] Processus de création ou de transformation de variables à partir des données brutes pour améliorer les performances des modèles.

  \item[\textbf{Feature importance}] Mesure de la contribution de chaque variable à la prédiction d'un modèle.

  \item[\textbf{Hyperparamètre}] Paramètre d'un modèle fixé avant l'entraînement (ex.~: $\lambda$ pour Ridge, $k$ pour K-Means).

  \item[\textbf{Overfitting (surapprentissage)}] Phénomène où un modèle apprend trop précisément les données d'entraînement et généralise mal sur de nouvelles données.

  \item[\textbf{Pipeline}] Enchaînement structuré d'étapes de traitement des données et de modélisation.

  \item[\textbf{Régularisation}] Technique pénalisant la complexité d'un modèle pour réduire l'overfitting (ex.~: Ridge $L_2$, Lasso $L_1$).

  \item[\textbf{SHAP}] Méthode d'interprétabilité basée sur la théorie des jeux coopératifs, attribuant à chaque variable une contribution marginale à la prédiction.

  \item[\textbf{Silhouette}] Coefficient mesurant la qualité d'un clustering~: proche de 1 si les clusters sont bien séparés, proche de 0 s'ils se chevauchent.

  \item[\textbf{SMOTE}] Technique de rééchantillonnage qui génère synthétiquement des exemples de la classe minoritaire pour corriger le déséquilibre des classes.

  \item[\textbf{Tuning}] Processus d'optimisation des hyperparamètres d'un modèle, souvent via une recherche aléatoire ou par grille.
\end{description}

\newpage
